% Ad Familiares 8.16 --- headnote
% ad-familiares-08-16, 16 April 49 BC, Intimilium (Liguria)

\textit{M. Caelius Rufus to Cicero, written from
Intimilium on the Ligurian coast on the sixteenth day
before the Kalends of May (16 April) 49 BC (manuscript
dateline \textit{Scr. Intimili xvi K. Mai. a. 705 (49)}).
This is the famous defection letter --- one of the
most-cited pieces of evidence in the entire correspondence
for how rapidly the political class of the late Republic
realigned in the spring of 49. Caelius is in Caesar's
camp, on the march up into Liguria and over the Alps en
route to the Spanish campaign that will decide the war
in the West. Cicero, whom Caelius last saw still
hesitating, has sent a letter that conveyed, without
saying so outright, that he had resolved to cross over
to Pompey. Caelius writes back the same day, ``stunned,''
to talk him out of it.}

\textit{The voice is the Caelian voice raised in alarm:
breezy when he needs to be (the gibe in section 2 that
Cicero is ``ashamed of being too little of an optimate''),
brutally direct when it counts (Caesar ``thinks --- and
even speaks --- of nothing but harshness and savagery'').
Caelius works every argument he has: family
(\textit{filius unicus}, the son-in-law Dolabella, the
house, the hopes); the political logic of the moment
(Caesar already takes Cicero's hesitation as offence;
to act now is to join the routed party after refusing to
follow them when they stood firm); the strategic forecast
(once Spain falls, what can the Pompeians offer?); and
his own dignity (had Caesar not been bringing him to
Spain, Caelius says, he would have ridden to Cicero
himself and held him back ``with all my force''). The
last section comes down to the compromise he wants Cicero
to accept --- not a defection to Caesar, but withdrawal
to some small town free of the war, which is the
position Cicero will in fact occupy for the rest of the
year, and which will protect him after Pharsalus. The
letter is also a striking piece of contemporary testimony
on the limits of Caesar's much-advertised clemency: the
Corfinium pardons did not, in Caelius's hearing, foretell
what Caesar meant to do after a final victory. ``He
thinks --- and even speaks --- of nothing but harshness
and savagery.''}
